\section*{Problem 7: Decoding II --- Linear Algebra (Level 7)}

\subsubsection*{Mathematical part}
\textbf{Recall:} Assume exactly $v$ errors occurred. Each error happened at some unknown field element $X_i$. Instead of searching directly for these positions, we construct another polynomial, the \textbf{Error Locator Polynomial} $\Lambda(x)$, whose roots are precisely $X_1^{-1}, X_2^{-1}, \ldots, X_v^{-1}$.

How do we determine $\Lambda(x)$? The syndromes satisfy a remarkable system of linear equations involving only the unknown coefficients $\Lambda_1, \dots, \Lambda_v$ of the locator polynomial:
\[
\begin{bmatrix}
S_0 & S_1 & \cdots & S_{v-1} \\
S_1 & S_2 & \cdots & S_{v} \\
\vdots & \vdots & \ddots & \vdots \\
S_{v-1} & S_{v} & \cdots & S_{2v-2}
\end{bmatrix}
\begin{bmatrix}
\Lambda_v \\ \Lambda_{v-1} \\ \vdots \\ \Lambda_1
\end{bmatrix}
=
\begin{bmatrix}
-S_v \\ -S_{v+1} \\ \vdots \\ -S_{2v-1}
\end{bmatrix}
\]
Finding $\Lambda(x)$ reduces to solving this Toeplitz linear system $Ax = b$. 

We do not know $v$ (the number of errors) in advance. The Peterson--Gorenstein--Zierler (PGZ) algorithm assumes the maximum possible errors $t$, builds the matrix, and attempts to solve it. If the matrix is singular (determinant zero), it assumes fewer errors occurred, decrements $t$, and tries again.

\paragraph{Worked Example}
If we hypothesize that $v=2$ errors occurred, the code constructs a $2 \times 2$ Toeplitz matrix from the sequence of syndromes $S_0, S_1, S_2, S_3$. The system of equations is exactly:
\[
\begin{bmatrix}
S_0 & S_1 \\
S_1 & S_2
\end{bmatrix}
\begin{bmatrix}
\Lambda_2 \\
\Lambda_1
\end{bmatrix}
=
\begin{bmatrix}
-S_2 \\
-S_3
\end{bmatrix}
\]
By solving this linear system, we find the coefficients $\Lambda_1$ and $\Lambda_2$. The Error Locator Polynomial is then formed as $\Lambda(x) = \Lambda_2 x^2 + \Lambda_1 x + 1$. 

\subsubsection*{Implementation part}
\textbf{Task:} Implement \lstinline|pgz_error_locator|.

\textbf{Details:} 
You must build the matrix $A$ and vector $b$ using the syndromes, then pass them to the provided \lstinline|solve_linear(A, b)| utility located in \lstinline|algebra_utils.py|. 
\begin{itemize}
    \item Start with $t = \text{len(syndromes)} // 2$. Loop downwards to $1$.
    \item Build $A$ (a list of lists) and $b$ (a list) for the current $t$.
    \item Use \lstinline|try/except| around \lstinline|solve_linear(A, b)|. If it raises a \lstinline|ValueError| (singular matrix), \lstinline|continue| to the next smaller $t$.
    \item If successful, construct and return the polynomial $\Lambda(x) = 1 + \Lambda_1 x + \dots + \Lambda_v x^v$. Remember that \lstinline|solve_linear| returns $[\Lambda_v, \dots, \Lambda_1]$. You must reverse this list and prepend $1$ (\lstinline|field.one|) to it before passing it to the \lstinline|Polynomial| constructor to satisfy the Low-to-High coefficient ordering.
    \item If the loop exhausts without success (every $t$ produced a singular matrix), return the trivial polynomial $\Lambda(x) = 1$, i.e.\ \lstinline|Polynomial([field.one])|. This constant polynomial has no roots, so the subsequent Chien Search will find zero error positions --- exactly the correct answer when no errors occurred.
\end{itemize}

\begin{center}
\renewcommand{\arraystretch}{1.5}
\begin{tabularx}{\textwidth}{@{} >{\bfseries}l X c @{}}
\toprule
Function & Arguments & Returns \\ \midrule
pgz\_error\_locator & \lstinline|syndromes: list, field: ExtensionField| & \lstinline|Polynomial| \\
solve\_linear & \lstinline|A: list[list], b: list| & \lstinline|list| \\
\bottomrule
\end{tabularx}
\end{center}

\subsubsection*{Experimentation part}
\textbf{Task:} Ensure the PGZ algorithm correctly handles any number of errors up to the theoretical limit $t$.
\textbf{UI Guide:} In the dashboard, introduce up to $t$ errors. The UI will display the Toeplitz matrix dynamically. If you introduce more than $t$ errors, you will see the system catastrophically fail, a fundamental limit of information theory.

\vspace{1cm}
